\subsection{نتایج تجربی}
ارزیابی عملکرد ماژول فشرده‌ساز \lr{RPCA Compressor} بر روی مدل‌های زبانی استاندارد در پایگاه داده بنچ‌مارک در جدول \ref{tab:benchmark_rpca_compressor} گزارش شده است:

\begin{table}[H]
\centering
\small
\caption{نتایج بنچ‌مارک فشرده‌ساز مبتنی بر تحلیل مؤلفه‌های اصلی مستحکم (\lr{RPCA Compressor})}
\label{tab:benchmark_rpca_compressor}
\begin{tabular}{lccccc}
\toprule
\textbf{مدل پایه \lr{(Base Model)}} & \textbf{دقت \lr{(Accuracy)}} & \textbf{شاخص \lr{F1}} & \textbf{حافظه \lr{(Peak VRAM)}} & \textbf{توان \lr{(Throughput)}} & \textbf{وضعیت} \\
\midrule
\lr{DistilBERT} & \lr{90.37\%} & \lr{0.9067} & \lr{311.19 MB} & \lr{138,163.61 tok/s} & \lr{PASS} \\
\lr{RoBERTa} & \lr{72.02\%} & \lr{0.7724} & \lr{301.72 MB} & \lr{92,321.44 tok/s} & \lr{PASS} \\
\lr{GPT-2} & \lr{60.21\%} & \lr{0.6825} & \lr{637.67 MB} & \lr{65,063.68 tok/s} & \lr{PASS} \\
\lr{Qwen} & -- & -- & -- & -- & \lr{SKIP} \\
\bottomrule
\end{tabular}
\end{table}

بررسی داده‌های عددی بنچ‌مارک نشان‌دهنده استحکام چشمگیر این ماژول در تمامی مدل‌های ارزیابی‌شده است:
\begin{itemize}
    \item در مدل \lr{DistilBERT}، دقت مدل از مقدار پایه \lr{90.48\%} به \lr{90.37\%} رسیده است که تنها بیانگر افت ناچیز \lr{0.11\%} بوده و هم‌زمان توان عملیاتی قابل توجه \lr{138,163.61 tok/s} را با وضعیت \lr{PASS} ثبت نموده است.
    \item در مدل \lr{RoBERTa}، دقت \lr{72.02\%} با مصرف حافظه بهینه \lr{301.72 MB} و توان عملیاتی \lr{92,321.44 tok/s} به دست آمده و افت دقت در محدوده مجاز کنترل شده است.
    \item در مدل زبانی اتورگرسیو \lr{GPT-2}، دقت \lr{60.21\%} و شاخص \lr{F1} برابر با \lr{0.6825} با توان گذردهی بیش از ۶۵ هزار توکن بر ثانیه احراز شده است.
    \item آزمایش روی مدل \lr{Qwen} به دلیل محدودیت منابع حافظه اختصاص‌یافته در اسکریپت یکپارچه در این نسخه در وضعیت عدم اجرا (\lr{SKIP}) قرار گرفته است.
\end{itemize}

\subsection{محدودیت‌ها و تحلیل شکست}
با وجود نتایج پایدار، تحلیل رفتار محاسباتی ماژول \lr{RPCA Compressor} مشخص‌کننده محدودیت‌های فنی زیر است:
\begin{itemize}
    \item \textbf{سربار زمانی چرخه‌های \lr{SVD} در فاز آفلاین:} اجرای چندباره عملگر \lr{SVD} در هر گام از تکرارهای \lr{ALM} برای لایه‌های حجیم (نظیر لایه‌های پیش‌خوران \lr{MLP}) زمان پیش‌پردازش کالیبراسیون را به صورت خطی با تعداد تکرارها افزایش می‌دهد.
    \item \textbf{حساسیت هایپرپارامتر آستانه‌گذاری $\mu$:} چنانچه پارامتر $\mu$ بیش از حد کوچک انتخاب شود، مؤلفه تنک $\mathbf{S}$ المان‌های غیرپرت را نیز جذب کرده و منجر به تضعیف اطلاعات مفید در مؤلفه کم‌رتبه $\mathbf{L}$ می‌گردد؛ در مقابل، مقادیر بسیار بزرگ $\mu$ کارایی جداسازی ناهنجاری‌ها را مختل می‌سازد.
\end{itemize}
